\subsection*{6. ``Differentiable theories'' (in the manner of Nash) and ``tame theories''}
The sixth section first continues the foundational program of \S5:
p.~37 proposes a dévissage of a proper tame map, simultaneous with a
finite family of closed subspaces, and introduces ``espaces globalement
modérés'' through complements in compact tame spaces. The specified
comparison of compactifications is the map
$g:X'^{\wedge}\longrightarrow X^{\wedge}$ restricting to the stated
bijection over $X$; p.~38 uses scheme compactifications in the examples
and extends the proposed framework to maps without the initial properness
requirement. There is an index slip on p.~37: the filtration is explicitly
decreasing, so its inclusion maps
$Y^i\hookrightarrow Y^{i-1}$ send every point of $Y^i$ into
$Y^{i-1}$. Consequently the printed ``open stratum''
$Y^i\setminus Y^{i-1}$ is empty. The readers preserve the exponent
$i-1$. After reporting what he had learned from Mebkhout, Grothendieck
states the ``principe des choix anodins'' on pp.~38--39: his criterion
examines the homotopy types of the spaces in which arbitrary construction
choices are made. This leads from isotopy categories, retaining the
components of isomorphism spaces, to the full homotopy types in the
``$\infty$-isotopique'' viewpoint. The surface examples return to the
anabelian cases and to Ladegaillerie's thesis from \S3; the discussion of
1-isotopy categories gives the corresponding bicategorical language.
At the end of p.~39 and the beginning of p.~40 he explains why he has
argued at such length for foundational work he does not intend to make
his own immediate priority. He then turns back to ideas from the late
1960s inspired by Nash, replacing the axioms for tame subsets by axioms
for rings of germs and the resulting smooth manifolds. The rounded
ring-system symbol on pp.~40--42 has an unresolved alphabet assignment;
the readers reproduce its checked contour and keep the doubt marks.
The stability requirements successively concern differential constructions,
the Weierstrass preparation condition and the footnote's henselian
formulation (p.~40), analytic continuation, and integration of Pfaff
systems (p.~41). The refusal to require algebras over all of
$\mathbb R$ leaves room for a smaller field of constants, pointwise
residue fields, and a hierarchy of ``fonctions élémentaires'' and
``constantes élémentaires''; these notions return to the fields discussed
in \S5, pp.~32--33. On p.~42 he judges the interesting questions about
the relation between a finitely generated theory and its field of
rationality ``out of reach'', a judgment explicitly withdrawn in note (1)
on p.~52, then explains how the older differentiable
ideas have become relevant again through the tame program. The final
proposal associates a tame theory with a differentiable one and studies
the differentiable structure off a closed nowhere dense tame subset,
with agreement of two such structures away from a further nowhere dense
subset. The closing sentence presents this as a setting for the more
precise comparison questions raised by the stratified dévissage of \S5.
(\href{https://grothendieck.umontpellier.fr/web/viewer.html?file=../119.pdf\#page=40}{Copy 1, pp.~37--42, beginning at site 39};
\href{https://grothendieck.umontpellier.fr/web/viewer.html?file=../119.pdf\#page=41}{principle of choices on pp.~38--39, beginning at site 40};
\href{https://grothendieck.umontpellier.fr/web/viewer.html?file=../119.pdf\#page=43}{Nash-inspired axioms on pp.~40--42, beginning at site 42};
\href{https://grothendieck.umontpellier.fr/web/viewer.html?file=../119.pdf\#page=45}{closing comparison on p.~42, site 44}.)
